\chapter{Code and Reproducibility}
\label{chap:code_appendix}

This appendix documents the implementation used to obtain the results presented
in this report. The code is written in Python and relies primarily on PyTorch
for tensor operations, automatic differentiation, neural-network modules and
GPU execution \citep{paszke2019pytorch}. The purpose of this appendix is not to
reproduce every source line, but to explain the responsibilities of the main
files, the data exchanged between them, and the choices that affect the
scientific interpretation of the results.

The paths below are relative to the repository root. The two main workflow
files retain the spelling used in the repository:
\begin{itemize}
\item \lstinline|WSGM_Foward_and_Generator.py|;
\item \lstinline|SGM_Foward.py|.
\end{itemize}

\section{Software organization}

The implementation is divided into executable workflows and reusable
libraries. The workflows assemble a complete experiment, whereas the library
modules implement the model, data and evaluation primitives.

{\footnotesize
\begin{longtable}{@{}p{0.38\textwidth}p{0.54\textwidth}@{}}
\toprule
\textbf{Component} & \textbf{Role in the pipeline} \\
\midrule
\endhead
\lstinline|data_mining_CEA.py| & Loads the CEA HDF5 simulations, splits them by simulation, applies preprocessing and augmentation, and exports the image datasets. \\
\lstinline|SGM_Foward.py| & Builds and trains the direct image-space VP-SGM, then saves its weights and configuration. \\
\lstinline|SGM_Generator.py| & Reloads a direct SGM from its configuration and checkpoint, samples in batches, and stores generated tensors. \\
\lstinline|WSGM_Foward_and_Generator.py| & Trains the coarse model and the conditional wavelet detail models, or runs the complete coarse-to-fine cascade according to \lstinline|WSGM_Config.json|. \\
\lstinline|PhyFID_Generator.py| & Trains the physical feature encoder and stores the validation statistics used by PhyFID. \\
\lstinline|lib/diffusion_lib/| & Generic DDPM/SGM models, U-Net, schedules, data loaders, training loops and experiment utilities. \\
\lstinline|lib/wavelet_diffusion_lib/| & Conditional diffusion models and training utilities specific to wavelet coefficients. \\
\lstinline|lib/PhyFID/| & Preprocessing, physical feature extraction and Frechet-distance evaluation. \\
\lstinline|Visualisation_fusionnee.ipynb| & Post-processing notebook used to compare generated ensembles and produce the report figures. \\
\bottomrule
\end{longtable}
}

The normal execution order is therefore
\[
\begin{aligned}
\text{CEA files} &\longrightarrow \text{processed tensors}
\longrightarrow \text{trained checkpoints} \\
&\longrightarrow \text{generated ensembles}
\longrightarrow \text{physical diagnostics and figures}.
\end{aligned}
\]
The saved configuration files are part of this chain: a checkpoint alone is
not sufficient to reconstruct a model if its architecture, normalization or
sampling parameters are not known.

\section{Dataset construction}

The entry point for dataset preparation is
\lstinline|data_mining_CEA.py|. It loads the three CEA HDF5 files through the
helpers in \lstinline|lib/cea_lib|. The function
\lstinline|split_indices_by_simulation| identifies a trajectory through the
fixed parameters stored in the label columns: Atwood number, phase and number
of modes. All snapshots sharing these parameters are assigned to the same
subset. This is important because a snapshot-level random split would put
different times from one physical trajectory in both training and validation,
and would make the validation error optimistic.

The resulting simulation-level split uses the ratios $0.8$, $0.1$ and $0.1$.
Each subset is then passed to \lstinline|data_preprocessing| and
\lstinline|data_augmentater| from \lstinline|lib/cea_lib|. The processed fields
are resized to $64\times64$ pixels. For the image-space diffusion workflow,
\lstinline|to_uint8| applies a per-field min--max rescaling before the tensors
are saved as \lstinline|training.pt|, \lstinline|validation.pt| and
\lstinline|test.pt|. The generic loader later maps these values to
$[-1,1]$, which is the convention expected by the image-space diffusion
models.

The wavelet workflow uses a different representation. Its coefficient files
are stored as floating-point tensors with shape $(N,4,H',W')$, where the four
channels are the approximation coefficient $cA$ and the detail coefficients
$cH$, $cV$ and $cD$. These coefficients must not be passed through the image
loader's uint8 normalization. The distinction is implemented by
\lstinline|WaveletApproximationDataset| and by the wavelet utilities described
below. The discrete wavelet transform is used as a multiresolution change of
representation, following the standard wavelet framework \citep{Mallat2009}
and its use for localized multiscale analysis in turbulence \citep{Farge1992}.

\section{Direct image-space SGM}

\subsection{Training workflow}

\lstinline|SGM_Foward.py| defines the direct image-space baseline. Its
\lstinline|Config| dataclass centralizes the dataset path, image shape,
training hyperparameters, VP-SDE schedule, U-Net architecture and output
paths. The function \lstinline|build_sgm| creates a time-conditioned U-Net and
wraps it in the \lstinline|SGM| class. \lstinline|main| then loads the
training and validation data, optionally restores an input checkpoint,
constructs an \lstinline|AdamW| optimizer, and calls
\lstinline|sgm_training_loop|. AdamW is the decoupled weight-decay optimizer
introduced by Loshchilov and Hutter \citep{loshchilov2019adamw}.

The continuous VP-SDE implemented in \lstinline|lib/diffusion_lib/SGM.py| is
defined by
\[
    \mathrm{d}x_t
    = -\frac{1}{2}\beta(t)x_t\,\mathrm{d}t
      + \sqrt{\beta(t)}\,\mathrm{d}W_t,
\]
with an affine schedule between \lstinline|beta_min| and
\lstinline|beta_max|. Because the marginal distribution is known in closed
form, the forward training step directly constructs
\[
    x_t = \sqrt{\bar{\alpha}(t)}x_0
          + \sqrt{1-\bar{\alpha}(t)}\,\epsilon,
    \qquad \epsilon\sim\mathcal{N}(0,I),
\]
without simulating the SDE with Euler--Maruyama. The score network is trained
with denoising score matching. The implementation supports both
\lstinline|epsilon|-prediction and the bounded $v$-parameterization; the
experiments reported in the later chapters use the configuration selected in
the corresponding JSON file.

The training loop in \lstinline|lib/diffusion_lib/training_loop.py| samples a
continuous time and Gaussian perturbation for every batch, evaluates an MSE
against the selected target, clips gradients when requested, evaluates the
validation loss, and saves the best checkpoint. This follows the score-based
diffusion formulation of Song et al. \citep{Song2021SDE}; the relation between
forward and reverse-time diffusion processes goes back to Anderson
\citep{ANDERSON1982313}. The older discrete implementation in
\lstinline|DDPM.py| follows the DDPM forward and reverse construction of Ho et
al. \citep{ho2020ddpm}; it remains available as a reusable library component,
although the main baseline discussed in this report is the continuous SGM.

\subsection{U-Net and time conditioning}

The model in \lstinline|lib/diffusion_lib/UNet.py| is an encoder--decoder
convolutional network with skip connections. At each encoder level, a stack of
convolutional blocks extracts features and a strided operation reduces the
spatial resolution. The bottleneck processes the coarsest representation. The
decoder upsamples the features and concatenates them with the corresponding
encoder activations. This contracting/expanding structure follows the U-Net
design of Ronneberger et al. \citep{ronneberger2015unet}, adapted here for
noise prediction rather than segmentation.

The diffusion time is represented by sinusoidal features and projected into
the channel dimension at each resolution level. In continuous-time mode, the
embedding receives $t\in[0,1]$ directly; in discrete DDPM mode, an embedding
table is used. The convolutional blocks can use layer normalization or group
normalization. The latter is motivated by the batch-size-independent
normalization proposed by Wu and He \citep{wu2018groupnorm}. The implementation
also exposes transposed-convolution or interpolation-based upsampling,
zero-padding or horizontal circular padding, and strided-convolution or
blur-pool downsampling.

The file \lstinline|lib/diffusion_lib/Attention.py| contains a spatial QKV
attention block based on scaled dot-product attention. The underlying
attention mechanism is associated with the Transformer formulation of Vaswani
et al. \citep{vaswani2017attention}, and the implementation follows the style
used in diffusion U-Nets such as the DDPM codebase \citep{ho2020ddpm}. There
is, however, an important implementation detail: the current
\lstinline|UNet.py| does not import or instantiate \lstinline|AttentionBlock|.
Consequently, the results reported in this thesis use the convolutional,
time-conditioned U-Net and do not use spatial attention. The attention file is
therefore documented for completeness, not as an active part of the trained
architecture.

\subsection{Generation workflow}

\lstinline|SGM_Generator.py| is deliberately separated from training. It
loads the JSON configuration generated by \lstinline|SGM_Foward.py|,
reconstructs the same U-Net and SGM objects, loads the saved state dictionary,
and switches the model to evaluation mode. \lstinline|generate_sgm_dataset|
generates samples in batches to limit memory usage. The SGM sampler integrates
the probability-flow ODE using either Euler or Heun; Heun is the default and
uses two score evaluations per step. The generated fields and a second dataset
sampled from a freshly initialized, untrained SGM (called the noise reference
in the script) are saved as PyTorch tensors for later analysis.

\section{Wavelet diffusion cascade}

\subsection{Conditional detail models}

The wavelet-specific implementation is in the directory
\lstinline|lib/wavelet_diffusion_lib|.

At each level, the coefficient tensor is split into a prior $cA$ and a detail tensor
$D=(cH,cV,cD)$. The class \lstinline|WaveletConditionalSGM| applies the VP-SDE
only to $D$. Its network receives the concatenation
\[
    [cA,D_t]
\]
and predicts the noise or $v$ target for the details. The approximation field
is therefore a condition and is not itself diffused in the conditional model.
Channel means and standard deviations are registered in the model so that
normalization used during training can be inverted before reconstruction.

The function \lstinline|wave_sgm_loss| implements conditional denoising score
matching. It samples $t$, adds the exact VP perturbation to the detail
channels, concatenates the unchanged prior, and computes the MSE between the
network output and the selected target. The corresponding training loop saves
the checkpoint with the lowest validation loss. The parallel
\lstinline|WaveletConditionalDDPM| and \lstinline|wave_training_loop| classes
provide the discrete DDPM alternative, with the same coarse/detail channel
convention.

\subsection{Coarse-to-fine orchestration}

\lstinline|WSGM_Foward_and_Generator.py| is the main orchestration script. Its
configuration distinguishes two model families:

\begin{description}
\item[Coarse SGM.] An unconditional image-space SGM is trained on the
  approximation coefficients $cA$. It supplies the initial coarse state for a
  complete generated cascade.
\item[Conditional detail SGM.] One model is trained at each wavelet level to
  generate details conditional on the current approximation coefficient.
\end{description}

The function \lstinline|generate_cascade| selects the initial coarse state,
generates the detail coefficients, applies the inverse discrete wavelet
transform, and passes the reconstructed approximation to the next level. The
two initial-state options are important for interpreting the experiments. In
the generated setting, $cA$ comes from the coarse diffusion model and the
cascade is fully generative. In the reference setting, $cA$ is read from a
validation field. The latter is an oracle-coarse ablation: it isolates the
conditional detail models and tests whether they can generate compatible fine
scales when the coarse state is already correct. It is not an independent
generative model and should not be counted as an end-to-end sampling method.

This factorization follows the Wavelet Score-Based Generative Model proposed
by Guth et al. \citep{Guth2022WSGM}. The WSGM paper motivates the use of
coarse-to-fine wavelet generation to reduce the cost of sampling, but the
present implementation and dataset determine whether that potential is
realized for RTI fields. In particular, the speed advantage and the quality of
the complete cascade remain separate empirical quantities.

\section{PhyFID implementation}

The files in \lstinline|lib/PhyFID| implement a domain-specific distribution
metric. \lstinline|preprocessing.py| converts input arrays to NCHW floating
point tensors in $[0,1]$. \lstinline|encoder.py| defines a convolutional
autoencoder. The encoder is trained on a training subset using reconstruction
MSE; after training, only its bottleneck vector is used as a physical feature.
\lstinline|metrics.py| computes the mean and covariance of these feature
vectors and evaluates the Frechet distance
\[
    d^2 = \lVert\mu_1-\mu_2\rVert_2^2
    + \operatorname{tr}\left(\Sigma_1+\Sigma_2
    -2(\Sigma_1\Sigma_2)^{1/2}\right).
\]
The construction is inspired by the Fréchet Inception Distance of Heusel et
al. \citep{heusel2017fid}, but it is not the standard FID: the feature space is
learned from RTI fields by the local autoencoder rather than taken from an
ImageNet Inception network. Consequently, PhyFID values are meaningful only
when the same encoder, preprocessing convention and comparable ensemble size
are used for all compared datasets.

\lstinline|PhyFID_Generator.py| creates this reference in four steps: it loads
the training and validation coefficient datasets, selects the requested
channel, trains the encoder, saves its checkpoint, and stores the validation
mean and covariance. The script contains this call at module level, so running
the file or importing it executes the reference-building procedure. Evaluation
of an already generated dataset is instead provided by
\lstinline|evaluate_phyfid_dataset| and does not retrain the encoder.

\section{Analysis notebook}

The notebook \lstinline|analyse_reslutat/Visualisation_fusionnee.ipynb| is an
analysis and figure-generation notebook rather than a second model implementation. Its
cells load saved tensors and metadata from the SGM and WSGM workflows, select
the cascade variant, compute lightweight statistical diagnostics, and export
the figures used in Chapter~\ref{chap:experiments}. Depending on the selected
configuration, the notebook compares generated and reference-coarse cascades,
plots coefficient or reconstruction fields, computes density statistics and
spectral diagnostics, and can call PhyFID.

For reproducibility, the notebook must be run against the checkpoints,
normalization statistics and generated tensors produced by the corresponding
experiment. Changing the cascade key, wavelet level, number of samples or
normalization changes the numerical comparison even if the plotting cells are
unchanged. The notebook should therefore be regarded as a post-processing
layer whose inputs are the serialized outputs of the scripts documented
above.

\section{Reproducibility checklist}

The following quantities should be preserved together with a reported result:

\begin{itemize}
\item the source simulation files and the simulation-level split seed;
\item the preprocessing, augmentation and per-channel normalization rules;
\item the JSON model configurations and the corresponding state dictionaries;
\item the initial coarse-state choice (reference or generated);
\item the wavelet family and decomposition levels;
\item the sampler, solver and number of reverse steps for every model;
\item the PhyFID encoder, reference statistics and number of evaluated samples.
\end{itemize}

These items are not merely software metadata. The initial coarse-state choice
changes the scientific question, and the reverse-step count changes the
measured speed--quality trade-off. In particular, a comparison between direct
SGM and WSGM is only an end-to-end comparison when the coarse coefficient is
also generated by the WSGM pipeline.
